% !TEX root = ../main.tex
\section{Numerical results}\label{sec:results}
In order to present our results, we consider benchmark scenarios with one or two non-vanishing RPV 
couplings. Additional non-zero RPV couplings could allow for further neutralino decay channels to 
open up; these would modify the relevant branching ratios and neutralino decay length, and hence, 
the presented sensitivity limits.



\subsection{One-coupling scenarios}
%
We first consider scenarios where only one non-zero RPV operator contributes to the relevant 
physical process at a time. For each of the RPV couplings in Eq.~\eqref{eq:W_RPV} (except the $\lam''
_{ijk}$'s), we list the relevant HNL searches providing constraints in~\cref{tab:single_couplings_BMs}.
We now discuss the sensitivity limits for each category of RPV coupling in detail.

\begin{table}[!ht]
\centering
\begin{adjustbox}{width=\columnwidth}
\begin{tabular}{|c|c|c|c|c|}
\hline
\bf{Coupling} & \bf{Direct Decays} & $\bm{E_T^\text{miss}}$ & \bf{DV} & \bf{Label}\\
\hline
$\kappa^i$ & $\nu$-mixing~\cite{Fernandez-Martinez:2023phj} & $\nu$-mixing~\cite{Fernandez-Martinez:2023phj} & $\nu$-mixing~\cite{Fernandez-Martinez:2023phj} & $\bm{\kappa}$\\
\hline
$\lambda'_{a11}$ & $\pi^\pm\to e^\pm_a + \tilde{\chi}^0_1$
& $\pi^0\to \text{invis.}$~\cite{NA62:2020pwi} & $\bm{\times}$ & \multirow{4}{*}{$\bm{\lambda'_{\pi}}$}\\[0.13 cm]
&~\cite{PiENu:2015seu, PIENU:2017wbj, PIENU:2019usb, PIONEER:2022alm} &&& \\[0.15 cm]
$\lambda'_{311}$ & $\left(\bm{\times}\right) \tau^\pm\to \{\pi^\pm/\rho^\pm\} + \tilde{\chi}^0_1$ & $\pi^0\to \text{invis.}$~\cite{NA62:2020pwi} & $\tau^\pm\to \pi^\pm + \tilde{\chi}^0_1; \tilde{\chi}^0_1 \to \nu_{\tau} + \{\pi^0/\rho^0/\eta/\eta'/\omega\}$
& \\[0.15 cm]
&&&\cite{Ballett:2019bgd, Kling:2018wct, DeVries:2020jbs, Beltran:2023nli}&\\
\hline
$\lambda'_{a12}$ & $K^\pm\to e^\pm_a + \tilde{\chi}^0_1$ & $K^0_L\to \text{invis.}$~\cite{ParticleDataGroup:2022pth} & $\bm{\times}$ & \multirow{4}{*}{$\bm{\lambda'_{K}}$}\\[0.15 cm]
&~\cite{Asano:1981he, Hayano:1982wu, NA62:2020mcv, NA62:2021bji} &&&\\[0.15 cm]
$\lambda'_{312}$ & $\left(\bm{\times}\right) \tau^\pm\to K^\pm + \tilde{\chi}^0_1$ & $K^0_L\to \text{invis.}$~\cite{ParticleDataGroup:2022pth} & $\bm{\times}$ & \\[0.15 cm]
\hline
$\lambda'_{i13}$ & $\left(\bm{\times}\right) B^\pm\to e^\pm_i + \tilde{\chi}^0_1$ & $B^0\to \text{invis.}$~\cite{BaBar:2012yut} & $\bm{\times}$ & $\bm{\lambda'_{B1}}$\\[0.15 cm]
\hline
$\lambda'_{i21}$ & $\left(\bm{\times}\right) D^\pm\to e^\pm_i + \tilde{\chi}^0_1$ & $K^0_L\to \text{invis.}$~\cite{ParticleDataGroup:2022pth} & $\bm{\times}$ & $\bm{\lambda'_{D/K}}$\\[0.15 cm]
\hline
$\lambda'_{i22}$ & $\left(\bm{\times}\right) D_s^\pm\to e^\pm_i + \tilde{\chi}^0_1$ & $\bm{\times}$ & $D_s^\pm\to e^\pm_i + \tilde{\chi}^0_1; \tilde{\chi}^0_1\to \nu_e + \{\phi/\eta/\eta'\}$ 
& $\bm{\lambda'_{D_s}}$\\[0.15 cm]
&&& \cite{Ballett:2019bgd, Kling:2018wct, DeVries:2020jbs, Beltran:2023nli} &\\[0.15cm]
\hline
$\lambda'_{i23}$ & $\left(\bm{\times}\right) B_c^\pm\to e^\pm_i + \tilde{\chi}^0_1$ & $B_s^0\to \text{invis.}$~\cite{ParticleDataGroup:2022pth} & $\bm{\times}$ & $\bm{\lambda'_{B_c/B_s}}$\\[0.15 cm]
\hline
$\lambda'_{i31}$ & $\bm{\times}$ & $B^0\to \text{invis.}$~\cite{BaBar:2012yut} & $\bm{\times}$ & $\bm{\lambda'_{B2}}$\\[0.15 cm]
\hline
$\lambda'_{i32}$ & $\bm{\times}$ & $B^0_s\to \text{invis.}$~\cite{ParticleDataGroup:2022pth} & $\bm{\times}$ & $\bm{\lambda'_{B_s}}$\\[0.15 cm]
\hline
$\lambda'_{i33}$ & $\bm{\times}$ & $\bm{\times}$ & $\bm{\times}$ & $\bm{-}$\\[0.15 cm]
\hline
$\lambda_{12a}$ & $\bm{\times}$ & $\mu^\pm \to e^\pm + \text{invis.}$
& $\bm{\times}$ & $\bm{\lambda_{\mu}}$\\[0.15 cm]
&&\cite{ParticleDataGroup:2022pth}&& \\
\hline
$\lambda_{123}; \lambda_{13i};$ & \multirow{2}{*}{$\bm{\times}$} & \multirow{2}{*}{$\tau^\pm \to \{e^\pm/\mu^\pm\} + \text{invis.}$}
& \multirow{2}{*}{$\bm{\times}$} & \multirow{2}{*}{$\bm{\lambda_{\tau}}$}\\[0.21 cm]
$\lambda_{232}; \lambda_{233}$ & & \cite{ParticleDataGroup:2022pth} & &\\[0.15 cm]
\hline
$\lambda_{231}$ & $\bm{\times}$ & $\{\tau^\pm/\mu^\pm\} \to e^\pm + \text{invis.}$
& $\bm{\times}$ & $\bm{\lambda_{\tau/\mu}}$\\[0.15 cm]
&&\cite{ParticleDataGroup:2022pth}&& \\
\hline
\end{tabular}
\end{adjustbox}
\caption{Details of the searches providing constraints when only one non-zero RPV operator contributes
at a time. We list all the bilinear, $LQ\Bar{D}$, and $LL\Bar{E}$ operators in the first column (by coupling).
The second to fourth columns contain the physical processes that provide constraints and the references to the relevant 
existing HNL searches targeting them. The fifth column indicates our labeling scheme for the scenarios.
$\bm{\times}$ denotes the absence of a constraining process, while $\left(\bm{\times}\right)$ labels that, 
in principle, the listed process may provide constraints but we could not find a relevant existing HNL search.
In the table, $a\in\{1,2\}$ and $i\in\{1,2,3\}$.}
\label{tab:single_couplings_BMs}
\end{table}


\subsubsection{Bilinear scenarios}
%
The bilinear couplings ($\kappa^i$ and the sneutrino vacuum expectation values) induce a mixing between 
the light bino and the three light neutrinos, \textit{cf}.~Eq.~\eqref{eq:nu_bino_bilinear_mixing}. Thus, 
all HNL searches constraining the mixing between the HNL and the neutrinos directly imply constraints on 
these couplings. The exclusion limits are shown in Fig.~\ref{fig:bilinear_bounds}. These have been read 
off from Ref.~\cite{Fernandez-Martinez:2023phj}. The bilinear couplings also generate mass terms for the 
neutrinos -- see, for instance, Refs.~\cite{Hall:1983id, Joshipura:1994ib, Nowakowski:1995dx, Banks:1995by, 
Allanach:2003eb} -- leading to the constraint $v_i,\,\kappa^i \lesssim \mathcal{O}\left(\SI{1}{\mega
\electronvolt}\right)$~\cite{Allanach:2003eb}; this is depicted as a gray horizontal line in the 
plot.\footnote{We note that the existing limit comes from the cosmological bound on the 
neutrino masses, and is thus scenario-dependent~\cite{ParticleDataGroup:2022pth}. Further, there is also dependence 
on undetermined supersymmetric parameters.} We find that the reinterpreted bounds from the HNL searches 
corresponding to the light-flavor neutrinos ($e$ and $\mu$) have sensitivity to regions beyond the current limits at low neutralino masses below 500 MeV.
In this region, the most constraining limits come from \texttt{PIENU}~\cite{PiENu:2015seu, PIENU:2017wbj,PIENU:2019usb}, \texttt{NA62}~\cite{NA62:2020mcv, NA62:2021bji}, and \texttt{T2K}~\cite{T2K:2019jwa, Arguelles:2021dqn}.
Beyond this mass, and for the tau case -- where charged kaon and pion decay into HNLs are kinematically forbidden -- existing limits are weaker.

\begin{figure}[!ht]
	\centering
	\includegraphics[width=0.8\textwidth]{figs/bilinear.pdf}
	\caption{Exclusion limits on bilinear RPV couplings as a function of the light bino mass, reinterpreted from existing HNL searches. The current limit on the bilinear couplings is shown as a horizontal gray line.}
	\label{fig:bilinear_bounds}
\end{figure}

\subsubsection{$LQ\bar{D}$ scenarios} 
%
For the remaining one-coupling scenarios, all existing limits on the RPV couplings are taken from 
Ref.~\cite{Dercks:2017lfq}. The $L_iQ_1\bar D_1$-operators couple to pions. With $i=1$, the decay $\pi^\pm 
\to e^\pm +\tilde{\chi}^0_1$ is allowed, for masses $m_{\tilde{\chi}^0_1} \le m_{\pi^{\pm}}-m_{e}$.\footnote{In
the small region of phase space, $m_{\pi^0}+m_{\nu_e} < m_{\tilde{\chi}^0_1} < m_\pi^\pm-m_e$, $\lam
'_{111}\not=0$ allows for the neutralino to decay into a neutrino and a pion; displaced-vertex 
searches can constrain this process. However, we have ignored this in~\cref{tab:single_couplings_BMs}
(also analogously for the kaon).} This process has been searched for in the context of HNLs, and the
most stringent existing limits are provided by \texttt{PIENU}~\cite{PIENU:2017wbj}. Additionally, 
the approved \texttt{PIONEER}~\cite{PIONEER:2022alm} experiment is projected to have sensitivity 
beyond these limits. We show the resulting contours in~\cref{fig:pion1}. The sharp drop in sensitivity 
at $m_{\tilde{\chi}^0_1} \approx \SI{65}{\MeV}$ occurs
because branching ratio measurements provide bounds below this mass, and peak searches above this threshold, \textit{cf.}~discussion in Sec.~\ref{sec:exp}. Similarly, for $i=2$, the pion can 
decay into a muon instead of an electron, in the mass range, $m_{\tilde{\chi}^0_1} \le m_{\pi^{\pm}}
-m_{\mu}$. This time, the sensitivity range, shown in~\cref{fig:pion2}, comes from 
\texttt{SIN}~\cite{Daum:1987bg} in addition to \texttt{PIENU}~\cite{PIENU:2019usb} and \texttt{PIONEER}.
Finally, with $i=3$, the decays, $\tau^{\pm} \to \pi^{\pm} (\rho^\pm) + \tilde{\chi}^0_1$, can occur if 
kinematically allowed. While we could not find a direct search for such a process involving an HNL, 
$\lam'_{311}\not=0$ additionally allows the neutralino to decay into a tau neutrino and one of 
$\pi^0/\rho^0/\eta/\eta'/\omega$. Thus, displaced-vertex searches at \texttt{DUNE}~\cite{Ballett:2019bgd}, 
\texttt{FASER(2)}~\cite{Kling:2018wct} and \texttt{MoEDAL-MAPP2}~\cite{DeVries:2020jbs,Beltran:2023nli} 
can show sensitivity.\footnote{See also Ref.~\cite{Dey:2020juy} for a proposed search at Belle~II~\cite{Belle-II:2010dht,Belle-II:2018jsg} for a light bino with non-vanishing RPV couplings $\lambda'_{311}$ or $\lambda'_{312}$.\label{footenote:belleII_LQD}}
We present the corresponding combined projected sensitivity reach in Fig.~\ref{fig:pion3}. In all the plots, we also show the best existing constraints on the RPV couplings for different sfermion masses.
For $\lam'_{111}$, this comes from neutrinoless double beta decay searches~\cite{Hirsch:1998kc,Bolton:2021hje}, and is rather stringent compared to the other two couplings.
Nevertheless, we see that the reinterpreted bounds easily outperform these existing constraints in the major part of the phase space region.
We also note that in all of the above scenarios, the RPV couplings contribute to the invisible decay width of the pion, via the decay into a neutrino and the bino, given the long lifetime of the latter. However, the limits obtained this way are not competitive compared with the above ones.

\begin{figure}[!ht]
\begin{subfigure}{0.49\textwidth}
  \includegraphics[width=\linewidth]{figs/single_coupling/pion_electronplusinv.pdf}
  \caption{$\lambda'_{111}$.}
  \label{fig:pion1}
  \end{subfigure}\hfill
\begin{subfigure}{0.49\textwidth}%
    \includegraphics[width=\linewidth]{figs/single_coupling/pion_muonplusinv.pdf}
  \caption{$\lambda'_{211}$.}
  \label{fig:pion2}
\end{subfigure}
\\
\centering
\begin{subfigure}{0.49\textwidth}
  \includegraphics[width=\linewidth]{figs/single_coupling/tau_1.pdf}
  \caption{$\lambda'_{311}$.
  }
  \label{fig:pion3}
  \end{subfigure}
\caption{Sensitivity limits on the $\bm{\lam'_{\pi}}$ one-coupling scenarios 
of~\cref{tab:single_couplings_BMs} as a function of the light bino mass, reinterpreted from 
existing HNL searches. Current (projected) limits obtained from the reinterpretation are 
shown as solid (dashed) gray lines. The existing limits on the RPV couplings are also shown 
in red, with the solid and dashed lines corresponding to varying assumptions of unknown SUSY mass scales.}
\label{fig:pion_single}
\end{figure}

Next, we consider the kaon scenarios, involving couplings of the type $\lambda'_{i12}$. For 
$i\in\{1,2\}$, as for the pion case above, the decay mode into an electron or a muon along 
with the bino opens up for the kaon, below the relevant kinematic thresholds. The most 
constraining current limits come from \texttt{KEK}~\cite{Asano:1981he, Hayano:1982wu} and 
\texttt{NA62}~\cite{NA62:2020mcv, NA62:2021bji}, and are depicted in Fig.~\ref{fig:kaon1} 
(electron case), and Fig.~\ref{fig:kaon2} (muon case).
The degradation at $m_{\tilde{\chi}^0_1}\sim \SI{200}{\MeV}$ in Fig.~\ref{fig:kaon2} occurs 
because the \texttt{KEK}-limit only applies up to this mass; beyond it the \texttt{NA62} 
limit applies. Once again, both $L_iQ_1\bar D_2$ also contribute to the invisible decays of the kaon but the 
resulting limits are weaker than those shown. However, for $\lam'_{312}$ 
-- where we could not find an existing HNL direct search for the charged decay, $\tau^\pm \to 
K^\pm + \tilde{\chi}^0_1$ (see Footnote~\ref{footenote:belleII_LQD}), the invisible width of $K^0_L$ can indeed be used to derive limits. 
Since there is no direct bound on this width, we use the uncertainty on the total measured 
width of $K^0_L$~\cite{ParticleDataGroup:2022pth} to estimate the upper bound, $\text{BR}(K^0_L \to 
\text{invis.}) < 4.1 \times 10^{-3}$. This, then, can be used to constrain $\lambda'_{312}$ 
since it induces the decay of $K^0_L$ into a neutrino and a bino. Analogously, the invisible 
width also provides limits on couplings of the type $\lambda'_{i21}$. Here also, we could 
not find existing direct searches for the relevant charged decay of $D$ into a lepton and 
an HNL. All these constraints are displayed in Fig.~\ref{fig:kaon3}. Once again, we can see 
that the reinterpreted bounds exclude regions of phase space orders of magnitude beyond 
those ruled out by the existing constraints on the RPV operators (also shown in the plots).

\begin{figure}[!ht]
\begin{subfigure}{0.49\textwidth}
  \includegraphics[width=\linewidth]{figs/single_coupling/kaon_electronplusinv.pdf}
  \caption{$\lambda'_{112}$.}
  \label{fig:kaon1}
  \end{subfigure}
  \begin{subfigure}{0.49\textwidth}
  \includegraphics[width=\linewidth]{figs/single_coupling/kaon_muonplusinv.pdf}
  \caption{$\lambda'_{212}$.}
  \label{fig:kaon2}
  \end{subfigure}
\\
\centering
\begin{subfigure}{0.49\textwidth}
  \includegraphics[width=\linewidth]{figs/single_coupling/kaon_intoinv.pdf}
  \caption{$\lambda'_{312}$ and $\lambda'_{i21}$.}
  \label{fig:kaon3}
  \end{subfigure}
\caption{As in~\cref{fig:pion_single} but for the $\bm{\lambda'_{K}}$ and $\bm{\lambda'_{D/K}}$ benchmarks of~\cref{tab:single_couplings_BMs}. The existing limits on the RPV couplings are shown in red, yellow, green, and blue.
}
\end{figure}
Similarly, for couplings of the type $\lambda'_{i13}$ and $\lambda'_{i31}$, we can use the 
invisible width of $B^0$. Here, direct measurements at \texttt{BaBar}~\cite{BaBar:2012yut} 
provide the stringent constraint, $\text{BR}(B^0 \to \text{invis.}) < 2.4\times 10^{-5}$. 
The resulting bounds are shown in Fig.~\ref{fig:B}.

\begin{figure}[!ht]
\begin{subfigure}{0.49\textwidth}
  \includegraphics[width=\linewidth]{figs/single_coupling/b_intoinv13.pdf}
  \caption{$\lambda'_{i13}$.}
  \label{fig:B1}
  \end{subfigure}
  \begin{subfigure}{0.49\textwidth}
  \includegraphics[width=\linewidth]{figs/single_coupling/b_intoinv_31.pdf}
  \caption{$\lambda'_{i31}$.
  }
  \label{fig:B2}
  \end{subfigure}
\caption{As in~\cref{fig:pion_single} but for the $\bm{\lambda'_{B1}}$ and $\bm{\lambda'_{B2}}$ benchmarks of~\cref{tab:single_couplings_BMs}. The existing limits on the RPV couplings are shown in red, green, and blue.
}
\label{fig:B}
\end{figure}

For the couplings $\lam'_{a22}$, with $a\in\{1,2\}$, we use potential future displaced-vertex 
searches from \texttt{DUNE}~\cite{Ballett:2019bgd}, \texttt{FASER2}~\cite{Kling:2018wct}, and 
\texttt{MoEDAL-MAPP2}~\cite{DeVries:2020jbs,Beltran:2023nli} since we could not find a direct 
search for the decay of $D^\pm_s$ into an HNL. In these scenarios, the neutralino decays via the 
RPV operator into a neutrino and one of $\phi/\eta/\eta'$. This implies the 
mass range, $\SI{548}{\MeV}\lesssim m_{\tilde{\chi}^0_1} < m_{D_s}-m_{e_a}$. Thus, the corresponding 
scenario with $\lam'_{322}\not=0$ involving a tau lepton is not possible. We show the resulting 
sensitivity limits in Fig.~\ref{fig:Ds}.

\begin{figure}[!ht]
\begin{subfigure}{0.48\textwidth}
  \includegraphics[width=\linewidth]{figs/single_coupling/D_11.pdf}
  \caption{$\lambda'_{122}$.}
  \label{fig:Ds1}
  \end{subfigure}\hfill
\begin{subfigure}{0.48\textwidth}%
    \includegraphics[width=\linewidth]{figs/single_coupling/D_12.pdf}
\caption{$\lambda'_{222}$.}
  \label{fig:Ds2}
\end{subfigure}
\caption{As in~\cref{fig:pion_single} but for the $\bm{\lambda'_{D_s}}$ benchmarks of~\cref{tab:single_couplings_BMs}.
}
\label{fig:Ds}
\end{figure}
For the remaining $LQ\Bar{D}$ operators, there are either no relevant processes providing constraints through meson or lepton decays ($\lambda'_{322}, \lambda'_{i33}$), or the obtained limits are not competitive with the existing bounds ($\lambda'_{i23}, \lambda'_{i32}$).


\subsubsection{$LL\Bar{E}$ scenarios} 
%
For the $\bm{\lam_{\mu}}, \bm{\lam_{\tau}}$ and $\bm{\lam_{\tau/\mu}}$ scenarios 
of~\cref{tab:single_couplings_BMs}, the corresponding operators 
contribute to the leptonic decays of the muon and the tau. Since we could not find direct 
measurements for the process $\mu^\pm\to e^\pm + \text{invis.}$ (and analogously for $\tau\to 
e$ and $\tau \to \mu$), where the invisible final state may be massive, we use the uncertainty
on the muon decay width~\cite{ParticleDataGroup:2022pth} to obtain an estimated bound, $\text{BR}(\mu^\pm
\to e^\pm +\text{invis.}) < 1.0\times 10^{-6}$. This leads to
constraints on the $\bm{\lam_{\mu}}$ and $\bm{\lam_{\tau/\mu}}$ scenarios, as shown in~\cref{fig:LLE}.
The analogous procedure with $\tau$ leads to limits weaker than the existing constraints and 
are hence not presented here.

\begin{figure}[!ht]
\centering	    
\includegraphics[width=0.6\textwidth]{figs/single_coupling/muon_intoinv.pdf}
\caption{As in~\cref{fig:pion_single} but for the $\bm{\lambda'_{\mu}}$ and $\bm{\lambda'_{\tau/\mu}}$ benchmarks of~\cref{tab:single_couplings_BMs}. The existing limits on the RPV couplings are shown in red, yellow, green, and blue.}
\label{fig:LLE}
\end{figure}







\subsection{Two-coupling scenarios}


Next, we consider scenarios where two RPV operators are simultaneously switched on; one corresponds to the production of the bino LSP, while the other leads to the decay. We 
can classify these scenarios based on the parent particle producing the neutralino; this then fixes 
the production RPV coupling. The relevant possibilities for the parent particles at beam-dump searches
and colliders are: pions, kaons, $D$ mesons, $D_s$ mesons, $B$ mesons, $B_c$ 
mesons, and $\tau$ leptons.\footnote{$\tau$ leptons are dominantly produced in the decay of the 
$D_s$ mesons at these experiments.} The corresponding production RPV couplings (and processes) 
can be read off from~\cref{tab:single_couplings_BMs}, \textit{e.g.}, $\lambda'_{a11}$ for the 
pion, etc.


The produced neutralino, owing to its long lifetime, then travels a certain macroscopic distance and 
decays via the other RPV coupling. Thus, displaced-vertex searches for HNLs are sensitive to such scenarios, if the final states match the decay products of 
the bino. We list, in~\cref{tab:decay_modes_HNL}, all relevant final states of such 
searches, along with the relevant kinematic thresholds at which the HNL can decay 
into them. Consulting the discussion in Sec.~\ref{sec:model}, we see that each of these final states 
can also arise from bino decays.

\begin{table}[!ht]
   \centering
\begin{tabular}{|lr|lr|}
\hline
Channel & Threshold & Channel & Threshold  \\
\hline
$\nu e^+ e^-$ & \SI{1.02}{\MeV} &$\mu^{\mp} K^{ \pm}$ & \SI{599}{\MeV} \\
$\nu e^{ \pm} \mu^{\mp}$ & \SI{105}{\MeV} & $\nu \rho^0$ & \SI{776}{\MeV} \\
$\nu \pi^0$ & \SI{135}{\MeV} & $e^{\mp} \rho^{ \pm}$ & \SI{776}{\MeV} \\
$e^{\mp} \pi^{ \pm}$ & \SI{140}{\MeV} & $\nu \omega$ & \SI{783}{\MeV} \\
$\nu \mu^{+} \mu^{-}$ & \SI{210}{\MeV} & $\mu^{\mp} \rho^{ \pm}$ & \SI{882}{\MeV} \\
$\mu^{\mp} \pi^{ \pm}$ & \SI{245}{\MeV} & $\nu \eta^{\prime}$ & \SI{958}{\MeV} \\
$e^{\mp} K^{ \pm}$ &  \SI{494}{\MeV} & $\nu \phi$ & \SI{1019}{\MeV} \\
$\nu \eta$ & \SI{548}{\MeV} &&\\
\hline
\end{tabular}
 \caption{Relevant final states from HNL (and bino) decay sorted by threshold mass. The active neutrino is considered massless.}
 \label{tab:decay_modes_HNL}
 \end{table}
In~\cref{tab:two_coupling_BMs}, we list, for each production category, the relevant RPV operator(s) leading to the final states of~\cref{tab:decay_modes_HNL}.
$\bm{\times}$ indicates that the given final state can not arise for the considered production mode, owing to kinematics.
This table can be used to identify all relevant two-coupling RPV scenarios that can be constrained by existing HNL DV searches. Note that we have not included bino production modes corresponding to the $LL\Bar{E}$ decays of the $\tau$ leptons in the table since these lead to weak limits, as in the one-coupling scenarios. We now discuss numerical results for a representative subset of the possibilities in the table. 

\subsubsection{Pion scenarios} 
%
For the $\pi^\pm$ category, we consider the two benchmark scenarios listed in~\cref{tab:pion_BMs}. 
The relevant mass range is identified by requiring that the production of the neutralino, and its 
subsequent decay (\textit{cf.}~\cref{tab:two_coupling_BMs}), should be both kinematically accessible.

We use this opportunity to explain one more subtlety of our recasting procedure for DV searches. In the HNL model, the relevant pion decay process occurs via the neutrino-mixings, 
\textit{e.g.}, $\pi^\pm \to e^\pm + N$ via $U_e$. However, a non-zero $U_e$ also induces 
the production of the HNL via decays of other particles, \textit{e.g.}, $K^\pm \to e^\pm 
+ N$ and the three-body decay, $K^\pm \to e^\pm + \pi^0 + N$. On the other hand, in our RPV 
benchmarks, this is not the case since $\lambda'_{i11}$ only couples to the pion. This affects
the kinematics of the HNL relative to the bino, and can be an issue for our simple scaling 
procedure, \textit{cf.} the discussion in Sec.~\ref{sec:exp}. For the pion benchmarks 
of~\cref{tab:pion_BMs}, this is not a problem since the kaon modes are sub-dominant $\left(
\mathcal{O}(1\%)\right)$ for the mass range identified above: $\SI{1}{\MeV} \lesssim m_N \lesssim 
\SI{139}{\MeV}$; see, for instance, Ref.~\cite{Berryman:2019dme} for a plot with the relevant 
branching ratios in the HNL scenario. However, in later benchmarks, we deal with this issue 
-- if it arises -- by restricting the mass range of the benchmark to ensure that the types of 
contributing parents are the same in both models (neglecting sub-dominant contributions up to 
$\mathcal{O}(10\%)$), so that the assumption of the same kinematics taken in the DV quick 
recasting method still holds. 

\begin{table}[H]
\centering
\begin{adjustbox}{width=0.75\columnwidth}
\rotatebox{90}{
\begin{tabular}{|c|c|c|c|c|c|c|c|c|c|c|c|}
\hline
\bf{Category} & \diagbox[width=5cm]{$\bm{\tilde{\chi}^0_1}$ \bf{Production}}{$\bm{\tilde{\chi}^0_1}$ \bf{Decay}} & $\bm{\nu ee}$ & $\bm{\nu e\mu}$ & $\bm{\nu \mu \mu}$ & $\bm{e\{\pi/\rho\}}$ & $\bm{\mu\{\pi/\rho\}}$ & $\bm{eK}$ & $\bm{\mu K}$ & $\bm{\nu\{\pi/\rho/\omega\}}$ & $\bm{\nu\{\eta/\eta'/\phi\}}$\\[0.15 cm]
\hline
%
\multirow{4}{*}{$\pi$} & $\lambda'_{111}: \pi^\pm \to e^\pm + \tilde{\chi}^0_1$ & \multirowcell{2}{$\lambda_{121}; \lambda_{131}$} & \multirowcell{2}{$\lambda_{121}; \lambda_{122}; \lambda_{321}; \lambda_{312}$} & \multirow{2}{*}{$\bm{\times}$} & \multirow{2}{*}{$\bm{\times}$} & \multirow{2}{*}{$\bm{\times}$} & \multirow{2}{*}{$\bm{\times}$} & \multirow{2}{*}{$\bm{\times}$} & \multirow{2}{*}{$\bm{\times}$} & \multirow{2}{*}{$\bm{\times}$}\\[0.15 cm]
& $(m_{\tilde{\chi}^0_1} \lesssim \SI{139}{\MeV})$ & & & & & & & & &\\[0.15 cm]
& $\lambda'_{211}: \pi^\pm \to \mu^\pm + \tilde{\chi}^0_1$ & \multirowcell{2}{$\lambda_{121}; \lambda_{131}$} & \multirowcell{2}{$\bm{\times}$} & \multirow{2}{*}{$\bm{\times}$} & \multirow{2}{*}{$\bm{\times}$} & \multirow{2}{*}{$\bm{\times}$} & \multirow{2}{*}{$\bm{\times}$} & \multirow{2}{*}{$\bm{\times}$} & \multirow{2}{*}{$\bm{\times}$} & \multirow{2}{*}{$\bm{\times}$}\\[0.15 cm]
& $(m_{\tilde{\chi}^0_1} \lesssim \SI{34}{\MeV})$ & & & & & & & & &\\[0.15 cm]
\hline
%
\multirow{4}{*}{$K$} & $\lambda'_{112}: K^\pm \to e^\pm + \tilde{\chi}^0_1$ & \multirowcell{2}{$\lambda_{121}; \lambda_{131}$} & \multirowcell{2}{$\lambda_{121}; \lambda_{122}; \lambda_{321}; \lambda_{312}$} & \multirow{2}{*}{$\lambda_{122}; \lambda_{322}$} & \multirow{2}{*}{$\lambda'_{111}$} & \multirow{2}{*}{$\lambda'_{211}$} & \multirow{2}{*}{$\bm{\times}$} & \multirow{2}{*}{$\bm{\times}$} & \multirow{2}{*}{$\lambda'_{i11}$} & \multirow{2}{*}{$\bm{\times}$}\\[0.15 cm]
& $(m_{\tilde{\chi}^0_1} \lesssim \SI{493}{\MeV})$ & & & & & & & & &\\[0.15 cm]
& $\lambda'_{212}: K^\pm \to \mu^\pm + \tilde{\chi}^0_1$ & \multirowcell{2}{$\lambda_{121}; \lambda_{131}$} & \multirowcell{2}{$\lambda_{121}; \lambda_{122}; \lambda_{321}; \lambda_{312}$} & \multirow{2}{*}{$\lambda_{122}; \lambda_{322}$} & \multirow{2}{*}{$\lambda'_{111}$} & \multirow{2}{*}{$\lambda'_{211}$} & \multirow{2}{*}{$\bm{\times}$} & \multirow{2}{*}{$\bm{\times}$} & \multirow{2}{*}{$\lambda'_{i11}$} & \multirow{2}{*}{$\bm{\times}$}\\[0.15 cm]
& $(m_{\tilde{\chi}^0_1} \lesssim \SI{388}{\MeV})$ & & & & & & & & &\\[0.15 cm]
\hline
\multirow{2}{*}{$D$} & $\lambda'_{a21}: D^\pm \to e_a^\pm + \tilde{\chi}^0_1$ & \multirowcell{2}{$\lambda_{121}; \lambda_{131}$} & \multirowcell{2}{$\lambda_{121}; \lambda_{122}; \lambda_{321}; \lambda_{312}$} & \multirow{2}{*}{$\lambda_{122}; \lambda_{322}$} & \multirow{2}{*}{$\lambda'_{111}$} & \multirow{2}{*}{$\lambda'_{211}$} & \multirow{2}{*}{$\lambda'_{112}$} & \multirow{2}{*}{$\lambda'_{212}$} & \multirow{2}{*}{$\lambda'_{i11}$} & \multirow{2}{*}{$\lambda'_{i11}; \lambda'_{i22}$}\\[0.15 cm]
& $(m_{\tilde{\chi}^0_1} \le m_{D}-m_{e_a})$ & & & & & & & & &\\[0.15 cm]
\hline
\multirow{2}{*}{$D_s$} & $\lambda'_{a22}: D_s^\pm \to e_a^\pm + \tilde{\chi}^0_1$ & \multirowcell{2}{$\lambda_{121}; \lambda_{131}$} & \multirowcell{2}{$\lambda_{121}; \lambda_{122}; \lambda_{321}; \lambda_{312}$} & \multirow{2}{*}{$\lambda_{122}; \lambda_{322}$} & \multirow{2}{*}{$\lambda'_{111}$} & \multirow{2}{*}{$\lambda'_{211}$} & \multirow{2}{*}{$\lambda'_{112}$} & \multirow{2}{*}{$\lambda'_{212}$} & \multirow{2}{*}{$\lambda'_{i11}$} & \multirow{2}{*}{$\lambda'_{i11}; \lambda'_{322}$}\\[0.15 cm]
& $(m_{\tilde{\chi}^0_1} \le m_{D_s}-m_{e_a})$ & & & & & & & & &\\[0.15 cm]
\hline
\multirow{2}{*}{$B$} & $\lambda'_{i13}: B^\pm(B^0) \to e_i^\pm(\nu_i) + \tilde{\chi}^0_1$ & \multirowcell{2}{$\lambda_{121}; \lambda_{131}$} & \multirowcell{2}{$\lambda_{121}; \lambda_{122}; \lambda_{321}; \lambda_{312}$} & \multirow{2}{*}{$\lambda_{122}; \lambda_{322}$} & \multirow{2}{*}{$\lambda'_{111}$} & \multirow{2}{*}{$\lambda'_{211}$} & \multirow{2}{*}{$\lambda'_{112}$} & \multirow{2}{*}{$\lambda'_{212}$} & \multirow{2}{*}{$\lambda'_{i11}$} & \multirow{2}{*}{$\lambda'_{i11}; \lambda'_{i22}$}\\[0.15 cm]
& $(m_{\tilde{\chi}^0_1} \le m_{B}-m_{e_i})$ & & & & & & & & &\\[0.15 cm]
\hline
\multirow{2}{*}{$B_c$} & $\lambda'_{i23}: B_c^\pm \to e_i^\pm + \tilde{\chi}^0_1$ & \multirowcell{2}{$\lambda_{121}; \lambda_{131}$} & \multirowcell{2}{$\lambda_{121}; \lambda_{122}; \lambda_{321}; \lambda_{312}$} & \multirow{2}{*}{$\lambda_{122}; \lambda_{322}$} & \multirow{2}{*}{$\lambda'_{111}$} & \multirow{2}{*}{$\lambda'_{211}$} & \multirow{2}{*}{$\lambda'_{112}$} & \multirow{2}{*}{$\lambda'_{212}$} & \multirow{2}{*}{$\lambda'_{i11}$} & \multirow{2}{*}{$\lambda'_{i11}; \lambda'_{i22}$}\\[0.15 cm]
& $(m_{\tilde{\chi}^0_1} \le m_{B_c}-m_{e_i})$ & & & & & & & & &\\[0.15 cm]
\hline
\multirow{4}{*}{$\tau$} & $\lambda'_{311}: \tau^\pm \to \{\pi^\pm/\rho^\pm\} + \tilde{\chi}^0_1$ & \multirowcell{2}{$\lambda_{121}; \lambda_{131}$} & \multirowcell{2}{$\lambda_{121}; \lambda_{122}; \lambda_{321}; \lambda_{312}$} & \multirow{2}{*}{$\lambda_{122}; \lambda_{322}$} & \multirow{2}{*}{$\lambda'_{111}$} & \multirow{2}{*}{$\lambda'_{211}$} & \multirow{2}{*}{$\lambda'_{112}$} & \multirow{2}{*}{$\lambda'_{212}$} & \multirow{2}{*}{$\lambda'_{a11}$} & \multirow{2}{*}{$\lambda'_{a11}; \lambda'_{i22}$}\\[0.15 cm]
& $(m_{\tilde{\chi}^0_1} \lesssim \SI{1637}{\MeV})$ & & & & & & & & &\\[0.15 cm]
& $\lambda'_{312}: \tau^\pm \to K^\pm + \tilde{\chi}^0_1$ & \multirowcell{2}{$\lambda_{121}; \lambda_{131}$} & \multirowcell{2}{$\lambda_{121}; \lambda_{122}; \lambda_{321}; \lambda_{312}$} & \multirow{2}{*}{$\lambda_{122}; \lambda_{322}$} & \multirow{2}{*}{$\lambda'_{111}$} & \multirow{2}{*}{$\lambda'_{211}$} & \multirow{2}{*}{$\lambda'_{112}$} & \multirow{2}{*}{$\lambda'_{212}$} & \multirow{2}{*}{$\lambda'_{i11}$} & \multirow{2}{*}{$\lambda'_{i11}; \lambda'_{i22}$}\\[0.15 cm]
& $(m_{\tilde{\chi}^0_1} \lesssim \SI{1283}{\MeV})$ & & & & & & & & &\\[0.15 cm]
\hline
\end{tabular}}
\end{adjustbox}
\caption{Relevant two-coupling RPV scenarios probed by HNL DV searches. Column one categorizes the parent meson/lepton; column two shows the 
corresponding RPV coupling, production process, and bino mass range; the remaining columns list the relevant final states of~\cref{tab:decay_modes_HNL} arising from bino decays and the corresponding decay RPV coupling(s) for each production category. $\bm{\times}$ indicates 
that the decay is kinematically disallowed.}
\label{tab:two_coupling_BMs}
\end{table}


\begin{table}
\centering
\begin{tabular}{|c|c|c|c|}
\hline
\bf{Label} & \bf{Production} & \bf{Decay} & $\bm{m_{\tilde{\chi}^0_1}}$\\
\hline
$\bm{\pi_1}$ & $\lambda'_{111}$ & $\lambda_{321}$ & $\SI{106}{\MeV} - \SI{139}{\MeV}$\\
$\bm{\pi_2}$ & $\lambda'_{211}$ & $\lambda_{121}$ & $\SI{1}{\MeV}-\SI{34}{\MeV}$\\
\hline
\end{tabular}
\caption{Details of the two-coupling RPV benchmark scenarios we study, corresponding to bino 
production from pions. The decay coupling in the third column leads to the final state that 
can be read off from~\cref{tab:two_coupling_BMs}. See the main text for details on how the 
mass range is determined.}
\label{tab:pion_BMs}
\end{table}



We show the reinterpreted limits for the benchmarks $\bm{\pi_1}$ and $\bm{\pi_2}$ in the RPV 
coupling vs.~mass plane in Fig.~\ref{fig:pion_two_1} and Fig.~\ref{fig:pion_two_2}, 
respectively. For two-dimensional visualization, we have set the production and decay 
couplings to be equal. The exclusion limits come from a corresponding HNL search at 
\texttt{Super-Kamiokande}~\cite{Coloma:2019htx}, as well as projections of sensitivity 
at \texttt{DUNE}~\cite{Ballett:2019bgd}. The former only constrains the mass range $m_N \gsim 
\SI{50}{\MeV}$ and, hence, does not have sensitivity to $\bm{\pi_2}$. We also show the 
existing bounds on the RPV couplings, taken from Ref.~\cite{Dercks:2017lfq}. There are no 
existing product bounds on the pairs considered \cite{Allanach:1999ic, Barbier:2004ez}. From 
Fig.~\ref{fig:pion_two_1}, it appears that the existing bound on $\lam'_{111}$ outperforms 
the reinterpreted bound. However, this is an artefact of the choice to set the production and decay couplings equal. We show the same exclusion limits again in 
Fig.~\ref{fig:pion_two_1a} -- this time in the coupling vs.~coupling plane for a fixed 
neutralino mass, $m_{\tilde{\chi}_1^0}$ = \SI{120}{\MeV}; one can see that the reinterpreted 
sensitivity projection can probe a small region of the phase space still allowed by the current 
limits.

\begin{figure}[!ht]
\begin{subfigure}{0.49\textwidth}
  \includegraphics[width=\linewidth]{figs/pion/pion_3.pdf}
  \caption{Limits in the RPV coupling vs.~bino mass plane for the benchmark $\bm{\pi_1}$ of~\cref{tab:pion_BMs}.}
  \label{fig:pion_two_1}
  \end{subfigure}\hfill
\begin{subfigure}{0.49\textwidth}%
    \includegraphics[width=\linewidth]{figs/pion/pion_2.pdf}
  \caption{Limits in the RPV coupling vs.~bino mass plane for the benchmark $\bm{\pi_2}$ of~\cref{tab:pion_BMs}.}
  \label{fig:pion_two_2}
\end{subfigure}
\\
\centering
\begin{subfigure}{0.49\textwidth}
  \includegraphics[width=\linewidth]{figs/pion/pion_3_two_coup.pdf}
  \caption{Limits in the production coupling vs.~decay coupling plane for the benchmark $\bm{\pi_1}$ of~\cref{tab:pion_BMs}, with neutralino mass fixed at \SI{120}{\MeV}.}
  \label{fig:pion_two_1a}
  \end{subfigure}
\caption{Current exclusion (solid lines) and projected sensitivity (dashed lines) limits corresponding to the two-coupling RPV scenarios with binos produced from pions; reinterpreted from HNL searches. The existing limits on the RPV couplings are shown in red, and green.}
\end{figure}





\subsubsection{Kaon scenarios}
%
Next, we study benchmarks corresponding to bino production via $K^\pm$ decays; the 
details are summarized in~\cref{tab:kaon_BMs}. Note, that for some benchmarks, 
(\textit{e.g.}, $\bm{K_3}$), the lower end of the mass range lies significantly 
above the kinematic threshold requirement of bino decay. This is, as discussed 
above, to ensure that kaons are the only parents in the HNL model, as they are in 
the RPV model.

\begin{table}[h]
\centering
\begin{tabular}{|c|c|c|c|}
\hline
\bf{Label} & \bf{Production} & \bf{Decay} & $\bm{m_{\tilde{\chi}^0_1}}$\\
\hline
$\bm{K_1}$ & $\lam'_{112}$ & $\lam'_{111}$ & $\SI{140}{\MeV} - \SI{493}{\MeV}$\\
$\bm{K_2}$ & $\lam'_{112}$ & $\lam'_{311}$ & $\SI{140}{\MeV}-\SI{493}{\MeV}$\\
$\bm{K_3}$ & $\lam'_{112}$ & $\lam_{321}$ & $\SI{140}{\MeV} - \SI{493}{\MeV}$\\
$\bm{K_4}$ & $\lam'_{212}$ & $\lam'_{211}$ & $\SI{140}{\MeV}-\SI{388}{\MeV}$\\
$\bm{K_5}$ & $\lam'_{212}$ & $\lam_{131}$ & $\SI{35}{\MeV} - \SI{388}{\MeV}$\\
\hline
\end{tabular}
\caption{As in~\cref{tab:pion_BMs} but for bino production from kaons.}
\label{tab:kaon_BMs}
\end{table}

The sensitivity limits for the kaon benchmarks are shown in Fig.~\ref{fig:kaon_two}. The 
single and -- wherever relevant -- product bounds (taken from Ref.~\cite{Allanach:1999ic, 
Barbier:2004ez}) on RPV couplings are also shown. Current exclusion limits are obtained 
by combining the results from existing HNL searches at \texttt{T2K}~\cite{T2K:2019jwa}, 
\texttt{Super-Kamiokande}~\cite{Coloma:2019htx}, \texttt{NuTeV}~\cite{NuTeV:1999kej}, 
and \texttt{MicroBooNE}~\cite{Kelly:2021xbv, MicroBooNE:2022ctm}, while the projections 
are all from \texttt{DUNE}~\cite{Ballett:2019bgd}.
In particular, for the benchmark $\bm{K_1}$, Ref.~\cite{Candia:2021bsl} has studied the sensitivity of \texttt{Super-Kamiokande} to the light binos in the RPV-SUSY, using existing data from the experiment.
Their results are found to be comparable with ours.
Once again, we see that the reinterpreted limits exclude (or are projected to probe) large swathes of 
parameter space allowed by the current bounds. The sharp reduction in sensitivity 
in~\cref{fig:kaon_two_5} below $m_{\tilde{\chi}^0_1} \approx \SI{150}{\MeV}$ 
arises because the most constraining current limit comes from \texttt{T2K} and only probes regions corresponding to $m_N \gsim \SI{150}{\MeV}$; below this the low-mass searches from \texttt{Super-K} and \texttt{MicroBooNE} provide exclusion.

\begin{figure}[!h]
\begin{subfigure}{0.49\textwidth}
  \includegraphics[width=\linewidth]{figs/kaon/kaon_1.pdf}
  \caption{Benchmark $\bm{K_1}$ from~\cref{tab:kaon_BMs}.}
  \label{fig:kaon_two_1}
  \end{subfigure}\hfill
\begin{subfigure}{0.49\textwidth}%
    \includegraphics[width=\linewidth]{figs/kaon/kaon_3.pdf}
  \caption{Benchmark $\bm{K_2}$ from~\cref{tab:kaon_BMs}.}
  \label{fig:kaon_two_2}
\end{subfigure}
\\
\begin{subfigure}{0.49\textwidth}
  \includegraphics[width=\linewidth]{figs/kaon/kaon_11.pdf}
  \caption{Benchmark $\bm{K_3}$ from~\cref{tab:kaon_BMs}.}
  \label{fig:kaon_two_3}
  \end{subfigure}\hfill
\begin{subfigure}{0.49\textwidth}%
    \includegraphics[width=\linewidth]{figs/kaon/kaon_5.pdf}
  \caption{Benchmark $\bm{K_4}$ from~\cref{tab:kaon_BMs}.}
  \label{fig:kaon_two_4}
\end{subfigure}
\\
\centering
\begin{subfigure}{0.49\textwidth}
  \includegraphics[width=\linewidth]{figs/kaon/kaon_8.pdf}
  \caption{Benchmark $\bm{K_5}$ from~\cref{tab:kaon_BMs}.}
  \label{fig:kaon_two_5}
  \end{subfigure}
\caption{Current exclusion (solid lines) and projected sensitivity (dashed lines) 
limits corresponding to the two-coupling RPV scenarios with binos produced from 
kaons in the RPV coupling vs.~bino mass plane; reinterpreted from HNL searches. The existing limits on the RPV couplings are shown in red and green (single bounds), and blue (product bound).
}
\label{fig:kaon_two}
\end{figure}







\subsubsection{$D$, ${D}_s$, and $\tau$ scenarios} 
%
\begin{figure}[h]
\begin{subfigure}{0.49\textwidth}
  \includegraphics[width=\linewidth]{figs/tau/tau_3.pdf}
  \caption{Benchmark $\bm{\tau_1}$ from~\cref{tab:d_BMs}.}
  \label{fig:tau_two_1}
  \end{subfigure}\hfill
\begin{subfigure}{0.49\textwidth}%
    \includegraphics[width=\linewidth]{figs/tau/tau_4.pdf}
  \caption{Benchmark $\bm{\tau_2}$ from~\cref{tab:d_BMs}.}
  \label{fig:tau_two_2}
\end{subfigure}
\caption{As in~\cref{fig:kaon_two} but for binos produced from $\tau$ leptons.}
\label{fig:tau_two}
\end{figure}
We summarize the details for the $\tau$ benchmarks we consider,
as well as the $D^\pm$ and $D^\pm_s$ meson ones in~\cref{tab:d_BMs}. We group them 
together in this section since, at the considered experiments, $\tau$ leptons are 
most copiously produced in the decays of the $D_s$ mesons. The corresponding 
sensitivity limits are shown in Fig.~\ref{fig:tau_two} for the $\tau$, and in Fig.~\ref{fig:D_two} for the mesons.
\begin{table}[h]
\centering
\begin{tabular}{|c|c|c|c|}
\hline
\bf{Label} & \bf{Production} & \bf{Decay} & $\bm{m_{\tilde{\chi}^0_1}}$\\
\hline
$\bm{\tau_1}$ & $\lambda'_{311}$ & $\lambda_{212}$ & $\SI{211}{\MeV} - \SI{1637}{\MeV}$\\
$\bm{\tau_2}$ & $\lambda'_{312}$ & $\lambda'_{311}$ & $\SI{140}{\MeV}-\SI{1283}{\MeV}$\\
$\bm{D_1}$ & $\lambda'_{122}$ & $\lambda'_{111}$ & $\SI{600}{\MeV} - \SI{1968}{\MeV}$\\
$\bm{D_2}$ & $\lambda'_{122}$ & $\lambda'_{211}$ & $\SI{600}{\MeV}-\SI{1968}{\MeV}$\\
$\bm{D_3}$ & $\lambda'_{122}$ & $\lambda'_{112}$ & $\SI{600}{\MeV} - \SI{1968}{\MeV}$\\
$\bm{D_4}$ & $\lambda'_{122}$ & $\lambda_{121}$ & $\SI{600}{\MeV}-\SI{1968}{\MeV}$\\
$\bm{D_5}$ & $\lambda'_{222}$ & $\lambda'_{211}$ & $\SI{600}{\MeV} - \SI{1863}{\MeV}$\\
$\bm{D_6}$ & $\lambda'_{222}$ & $\lambda_{131}$ & $\SI{600}{\MeV} - \SI{1863}{\MeV}$\\
$\bm{D_7}$ & $\lambda'_{221}$ & $\lambda_{232}$ & $\SI{260}{\MeV} - \SI{1764}{\MeV}$\\
\hline
\end{tabular}
\caption{As in~\cref{tab:pion_BMs} but for bino production from $\tau$ leptons, and $D$ and $D_s$ mesons.}
\label{tab:d_BMs}
\end{table}

For the $\tau$ lepton scenarios, the exclusion and projected search sensitivity, shown in~\cref{fig:tau_two}, come from \texttt{BEBC}~\cite{WA66:1985mfx, 
Barouki:2022bkt}, \texttt{CHARM}~\cite{CHARM:1985nku, Boiarska:2021yho}, and 
\texttt{ArgoNeuT}~\cite{ArgoNeuT:2021clc}; and \texttt{DUNE}~\cite{Ballett:2019bgd}, 
\texttt{FASER2}~\cite{Kling:2018wct,Beltran:2023nli}, and 
\texttt{MoEDAL-MAPP2}~\cite{DeVries:2020jbs, Beltran:2023nli}, respectively. There is 
no current search targeting the final state of $\bm{\tau_2}$, while the current limit 
on $\bm{\tau_1}$ beats even the projected search sensitivity at 
\texttt{DUNE} and \texttt{FASER}. The sharp drop in sensitivity in~\cref{fig:tau_two_1} 
at $m_{\tilde{\chi}_1^0} \approx \SI{1000}{\MeV}$ occurs because the $\rho$ and bino 
decay mode of the $\tau$ lepton in the RPV model (\textit{cf.}~\cref{tab:single_couplings_BMs}) 
becomes kinematically inaccessible, leading to the reduction in production.

\begin{figure}[h!]
\begin{subfigure}{0.45\textwidth}
  \includegraphics[width=\linewidth]{figs/D/D_1.pdf}
  \caption{Benchmark $\bm{D_1}$ from~\cref{tab:d_BMs}.}
  \label{fig:D_two_1}
  \end{subfigure}\hfill
\begin{subfigure}{0.45\textwidth}
   \includegraphics[width=\linewidth]{figs/D/D_3.pdf}
  \caption{Benchmark $\bm{D_2}$ from~\cref{tab:d_BMs}.}
   \label{fig:D_two_2}
\end{subfigure}\hfill
\\
\begin{subfigure}{0.45\textwidth}%
    \includegraphics[width=\linewidth]{figs/D/D_13.pdf}
  \caption{Benchmark $\bm{D_3}$ from~\cref{tab:d_BMs}.}
   \label{fig:D_two_3}
\end{subfigure}\hfill
\begin{subfigure}{0.45\textwidth}
  \includegraphics[width=\linewidth]{figs/D/D_29.pdf}
  \caption{Benchmark $\bm{D_4}$ from~\cref{tab:d_BMs}.}
   \label{fig:D_two_4}
  \end{subfigure}\hfill
  \\
\begin{subfigure}{0.45\textwidth}
   \includegraphics[width=\linewidth]{figs/D/D_8.pdf}
  \caption{Benchmark $\bm{D_5}$ from~\cref{tab:d_BMs}.}
   \label{fig:D_two_5}
\end{subfigure}\hfill
\begin{subfigure}{0.45\textwidth}%
    \includegraphics[width=\linewidth]{figs/D/D_19.pdf}
  \caption{Benchmark $\bm{D_6}$ from~\cref{tab:d_BMs}.}
   \label{fig:D_two_6}
\end{subfigure}
\\
\centering
\begin{subfigure}{0.45\textwidth}
   \includegraphics[width=\linewidth]{figs/D/D_24.pdf}
   \caption{Benchmark $\bm{D_7}$ from~\cref{tab:d_BMs}.}
    \label{fig:D_two_7}
\end{subfigure}
  \caption{As in~\cref{fig:kaon_two} but for binos produced from $D$ and $D_s$ mesons.}
  \label{fig:D_two}
\end{figure}

For the $D$ and $D_s$ mesons, current exclusions are provided by searches at 
\texttt{BEBC}~\cite{WA66:1985mfx, Barouki:2022bkt}, \texttt{CHARM}~\cite{CHARM:1985nku}, 
and \texttt{NuTeV}~\cite{NuTeV:1999kej}.
Further, \texttt{DUNE}~\cite{Ballett:2019bgd}, \texttt{FASER2}~\cite{Kling:2018wct}, and \texttt{MoEDAL-MAPP2}~\cite{DeVries:2020jbs,Beltran:2023nli} are projected to improve this reach. For all $D_s$ benchmarks except 
$\bm{D_7}$, we only consider neutralinos with mass, $m_{\tilde{\chi}^0_1} > \SI{600}{\MeV}$.
This is because either the above experiments only constrain the corresponding parameter 
region in the HNL model, or because kaons and pions dominate the HNL production for 
lower masses. One exception is \texttt{BEBC}, where nearly all produced pions and kaons 
are absorbed by a high-density target before they can decay, and HNL production is 
dominated by $D$ meson decays. Since benchmark $\bm{D_7}$ involves production from 
$D$-mesons, we can probe lower bino masses in our scenario.
The kink in the corresponding sensitivity limit in~\cref{fig:D_two_7} at $m_{\tilde{\chi}_1^0} \approx \SI{500}{\MeV}$ occurs because \texttt{CHARM} takes over from \texttt{BEBC}.
We note that the final states of benchmarks $\bm{D_1}$ and $\bm{D_3}$ are not covered by the existing searches but will be covered by the upcoming experiments. 


Once again, we see -- for all the plots in this section -- that reinterpreting existing 
and projected limits on HNL models in terms of our RPV scenarios gives bounds on the 
parameter space that improve upon existing limits by orders of magnitude.


\subsubsection{$B$ and $B_c$ scenarios} 
%
Here, the bino is produced in $B (B^\pm/\overset{\scriptscriptstyle(-)}{B^0})$ or $B^\pm_c$ 
decays, \textit{cf.}~\cref{tab:b_BMs}. While HNL production in these 
modes only becomes dominant above the $D$-meson thresholds, the projected search 
sensitivity for \texttt{FASER(2)} provided in Ref.~\cite{Kling:2018wct} shows the results separately for the HNLs from $B$-meson decays and from $D$-meson and kaon 
decays, enabling us to choose benchmarks with masses lower than the $D$ thresholds. Up 
to about $m_N \approx \SI{2700}{\MeV}$, the production is dominated by $B$ decays and we choose 
the first three benchmarks accordingly. Beyond this, $B_c$ decays also become 
significant, and for $m_N \gsim \SI{3500}{\MeV}$, they become the dominant modes; 
the last two benchmarks focus on this.

\begin{table}[h]
\centering
\begin{tabular}{|c|c|c|c|}
\hline
\bf{Label} & \bf{Production} & \bf{Decay} & $\bm{m_{\tilde{\chi}^0_1}}$\\
\hline
$\bm{B_1}$ & $\lambda'_{113}$ & $\lambda'_{122}$ & $\SI{548}{\MeV} - \SI{2700}{\MeV}$\\
$\bm{B_2}$ & $\lambda'_{113}$ & $\lambda_{131}$ & $\SI{160}{\MeV}-\SI{2700}{\MeV}$\\
$\bm{B_3}$ & $\lambda'_{213}$ & $\lambda'_{211}$ & $\SI{160}{\MeV} - \SI{2700}{\MeV}$\\
$\bm{B_4}$ & $\lambda'_{123}$ & $\lambda'_{311}$ & $\SI{3500}{\MeV}-\SI{6275}{\MeV}$\\
$\bm{B_5}$ & $\lambda'_{123}$ & $\lambda_{131}$ & $\SI{3500}{\MeV} - \SI{6275}{\MeV}$\\
\hline
\end{tabular}
\caption{As in~\cref{tab:pion_BMs} but for bino production from $B$ and $B_c$ mesons.}
\label{tab:b_BMs}
\end{table}

\begin{figure}[h]
\begin{subfigure}{0.48\textwidth}
  \includegraphics[width=\linewidth]{figs/B/B_31.pdf}
  \caption{Benchmark $\bm{B_1}$ from~\cref{tab:b_BMs}.}
  \end{subfigure}\hfill
\begin{subfigure}{0.48\textwidth}%
    \includegraphics[width=\linewidth]{figs/B/B_19.pdf}
  \caption{Benchmark $\bm{B_2}$ from~\cref{tab:b_BMs}.}
\end{subfigure}
\\
\begin{subfigure}{0.48\textwidth}
  \includegraphics[width=\linewidth]{figs/B/B_9.pdf}
  \caption{Benchmark $\bm{B_3}$ from~\cref{tab:b_BMs}.}
  \end{subfigure}\hfill
\begin{subfigure}{0.48\textwidth}%
    \includegraphics[width=\linewidth]{figs/B/B_6.pdf}
  \caption{Benchmark $\bm{B_4}$ from~\cref{tab:b_BMs}.}
\end{subfigure}
\\
\centering
\begin{subfigure}{0.48\textwidth}
  \includegraphics[width=\linewidth]{figs/B/B_20.pdf}
  \caption{Benchmark $\bm{B_5}$ from~\cref{tab:b_BMs}.}
  \end{subfigure}
\caption{As in~\cref{fig:kaon_two} but for binos produced from $B$ and $B_c$ mesons.}
\label{fig:B_two}
\end{figure}

The corresponding sensitivity limits are presented in Fig.~\ref{fig:B_two}.
There are no existing constraints; however, projections from \texttt{FASER2}~\cite{Kling:2018wct,Beltran:2023nli} and \texttt{MoEDAL-MAPP2}~\cite{DeVries:2020jbs,Beltran:2023nli} show that we should be able to probe the RPV parameter space up to 2-3 orders of magnitude beyond what is ruled out by current limits.